\appendix

\section{The log-likelihood}\label{app:likelihood}
Let $x(t)$ be the state at time $t$: the set of links, and the impact $I_a(t)$ of every existing action $a$.
Between events, impacts decay deterministically, $I_a(t) = I_a(t_0)\,e^{-\delta_a (t - t_0)}$, and an action ceases to exist when its impact reaches its threshold, at a deterministic time.
Forced events happen at given times, and triggered acts happen together with the contact events that trigger them.
All other events are random with rates $\lambda_e(x(t))$ that depend on the links and, through angles weighted by impacts, continuously on time.
Let $\Lambda(t)$ be their total rate.

\begin{proposition}\label{prop:likelihood}
Given the initial state and the forced events, the probability density of a sequence of random events $e_1, \dots, e_m$ at times $t_1 < \dots < t_m$ in $[0, T]$ is
\begin{equation*}
  p = \prod_{i=1}^m \lambda_{e_i}(x(t_i^-)) \, \exp\Bigl(-\int_0^T \Lambda(t)\,dt\Bigr),
\end{equation*}
and its logarithm is Eq.~\eqref{eq:logl}.
\end{proposition}

\begin{proof}
Between the times of random events, forced events, and expiries of actions, the state changes only deterministically, so the process is a point process with the stochastic intensity $\Lambda(t)$ with respect to its history, and an event at time $t$ is $e$ with probability $\lambda_e(x(t^-))/\Lambda(t)$.
The density of a point process with intensity $\Lambda$ is $\prod_i \Lambda(t_i) \exp(-\int_0^T \Lambda\,dt)$, and multiplying by the probabilities of the identities gives $p$.
Forced events and expiries happen at times determined by the past and the given input, and triggered acts with probability one, so they add no factors.
\end{proof}

\noindent \tricl\ computes the integral event by event: the rate of an event changes deterministically between its updates, and when an event is updated or removed, the integral of its rate since its last update is added by Gauss--Legendre quadrature (eight nodes on pieces of length at most $1/\delta_{\max}$, where $\delta_{\max}$ is the largest decay rate of its weighted angles, which gives a relative accuracy of about $10^{-13}$).
The gradient is accumulated in the same way.
For the symmetric relationship type of contacts, the two directions $e = (i, j)$ and $\bar e = (j, i)$ of a pair are two possible events that cause the same change; Proposition~\ref{prop:likelihood} is the likelihood of the sequence with the recorded direction of each event, and the likelihood without directions replaces $\lambda_{e_i}$ by $\lambda_{e_i} + \lambda_{\bar e_i}$.
Since all models in this paper have $\lambda_e = \lambda_{\bar e}$ in every state, the two log-likelihoods differ by the constant $m \ln 2$, and estimates, likelihood differences and standard errors are the same.
The competing model C1 is implemented with the same convention.

\section{Exactness of the simulation}\label{app:thinning}
The sum tree of rates holds, for every individually scheduled event $e$, a weight $w_e \ge \lambda_e(x(s))$ for all times $s$ until its next update: for constant rates $w_e = \lambda_e$, and for rates with decaying impacts $w_e = \bar A f(\bar B)$, with $\bar A$ the attempt rate at the last update (attempt coefficients are non-negative, so the attempt rate decreases as impacts decay) and $\bar B$ the probability units at the last update with all negative contributions of decaying angles omitted.
Pairs without adjacent angles are covered by one entry per event type, whose weight is an upper bound of the sum of their rates; an attempt of this entry draws a pair through a channel chosen proportionally to its attempt rate (a uniformly random pair, or a uniformly random link of a leg type and a uniformly random partner) and is accepted with the ratio of the pair's success probability to an upper bound.
Let $W$ be the total weight.

\begin{proposition}\label{prop:thinning}
If an attempt waits an exponential time with rate $W$, chooses an entry proportionally to its weight, performs a chosen individually scheduled event $e$ with probability $\lambda_e(x(t))/w_e$ (and otherwise refreshes its weight at time $t$), and processes forced events and expiries when their times are reached before the next attempt, then the produced events have the law of the process of Proposition~\ref{prop:likelihood}.
\end{proposition}

\begin{proof}
Between updates, the attempts form a Poisson process of rate $W$ whose points are independently marked; an attempt at time $t$ produces event $e$ with probability $\lambda_e(x(t))/W$, both for individually scheduled events and for covered pairs, since the channel probabilities and acceptance ratios multiply to $\lambda_e/W$.
By the thinning theorem \citep{lewis1979,ogata1981}, the produced events form a point process with intensity $\lambda_e(x(t))$ for each $e$, as long as $w_e \ge \lambda_e(x(s))$ on the interval, which holds by construction.
Refreshing a weight after a rejected attempt and processing deterministic events changes the bounds and the state only at stopping times, after which the argument applies anew with fresh exponential waiting times (the waiting time is memoryless).
\end{proof}

\section{Time-rescaling residuals}\label{app:residuals}
Under a correct model, the termination of a contact between $i$ and $j$ that started at time $s$ has the hazard $h(t) = \lambda_{(i,j)}(x(t)) + \lambda_{(j,i)}(x(t))$ given the path, so that $\Lambda = \int_s^{s+\tau} h(t)\,dt$, with $\tau$ the duration, is a unit exponential variable, and these variables are independent across contacts \citep{brown2002}.
\tricl\ accumulates the integrated rate of every termination event from the time it is added (when the contact starts) to its occurrence, including time-varying rates, and adds the integral of the mirrored event of the same pair.
Contacts still in progress at the end of a day are not included, which slightly favors short durations; the effect is small, since few contacts are in progress at 17:00.

\section{Pair closure}\label{app:closure}
For a class with $N$ students and $M = N(N-1)/2$ pairs, let $L_c$ be the number of pairs with marks $c$ in contact (out of $M_c$), $L = \sum_c L_c$, $U = M - L$, and $Z$ and $O$ the sums of the numbers $n$ of common contacts over pairs in contact and not in contact.
Since every two-path is counted once in $Z + O$, $\sum_i d_i (d_i - 1) = 2(Z + O)$ for the numbers of contacts $d_i$.
The closure assumes that, for a pair in contact, $n \sim \mathrm{Poisson}(Z/L)$ and the other contacts of the two students beyond the common ones are Poisson with mean $2 b_1$, $b_1 = (Z + O)/L - Z/L$ (the mean number of other contacts of an endpoint of a contact, $\sum_i d_i(d_i - 1)/\sum_i d_i$, minus the common ones), and for a pair not in contact $n \sim \mathrm{Poisson}(O/U)$ with $b_0 = [\sum_i d_i (N - 1 - d_i)]/(2U) - O/U$.
With $\lambda^{\pm}_c(n, m)$ the pair rates of formation and termination (both directions) for a pair with marks $c$, $n$ common contacts and $m$ other contacts of its endpoints, and $\langle \cdot \rangle_0$, $\langle \cdot \rangle_1$ the expectations under these distributions,
\begin{equation*}
\begin{aligned}
  \dot L_c &= (M_c - L_c) \langle \lambda^+_c \rangle_0 - L_c \langle \lambda^-_c \rangle_1, \\
  \dot Z &= \textstyle\sum_c \bigl[(M_c - L_c) \langle 3 n \lambda^+_c \rangle_0 - L_c \langle 3 n \lambda^-_c \rangle_1\bigr], \\
  \dot O &= \textstyle\sum_c \bigl[(M_c - L_c) \langle (m - n) \lambda^+_c \rangle_0 - L_c \langle (m - n) \lambda^-_c \rangle_1\bigr],
\end{aligned}
\end{equation*}
since a new contact with $n$ common contacts turns its own $n$ open two-paths into closed ones and adds two closed two-paths for each common contact ($\Delta Z = 3n$), and adds one open two-path for each other contact of its endpoints ($\Delta O = m - n$), and the termination of a contact reverses these changes.
The simple closure replaces the conditional distributions by the same Poisson distributions for both states, with means $(N - 2)\rho^2$ and $(N - 2)\rho$, $\rho = L/M$.
